\chapter{Glossário}
\label{ch:glossario}
\addcontentsline{toc}{chapter}{Glossário}

\begin{longtable}{@{}p{0.27\linewidth} p{0.68\linewidth}@{}}
\toprule
\textbf{Termo} & \textbf{Definição} \\
\midrule
\endfirsthead
\toprule
\textbf{Termo} & \textbf{Definição} \\
\midrule
\endhead

AHP & \emph{Analytic Hierarchy Process}. Método de priorização
proposto por \citet{saaty1980}, baseado em comparações aos pares
ponderadas por níveis hierárquicos de critérios; tem complexidade
$O(n^2)$ no número de itens. \\

\emph{AST} (\emph{Abstract Syntax Tree}) & Representação em árvore
da estrutura sintática de um programa, usada como ponto de partida
para análise estática de código-fonte. \\

\emph{Call Graph} & Grafo direcionado em que cada vértice é uma
função (ou método) e cada aresta $(f_i, f_j)$ indica que $f_i$
invoca $f_j$. \\

Cadeia de Markov & Processo estocástico em que a probabilidade do
estado seguinte depende apenas do estado atual. Subjacente aos três
algoritmos analisados nesta dissertação. \\

\emph{Dangling Node} & Vértice sem arestas de saída. No PageRank é
tratado por redistribuição uniforme da sua massa de probabilidade,
garantindo ergodicidade da cadeia de Markov. \\

Distribuição estacionária & Vetor probabilístico $\mathbf{r}^{*}$
que verifica $\mathbf{r}^{*} = M\mathbf{r}^{*}$, isto é, é
invariante sob aplicação da matriz de transição $M$. Existe e é
única em cadeias de Markov ergódicas. \\

Efeito TKC & \emph{Tightly Knit Community Effect}: vulnerabilidade
do HITS em que sub-comunidades densas concentram
desproporcionalmente os \emph{scores}; o SALSA foi projetado para
mitigá-lo. \\

Ergodicidade & Propriedade de uma cadeia de Markov que garante a
existência de uma distribuição estacionária única, independente do
estado inicial. \\

\emph{Feature module} & Agregação de funções e métodos do código
relacionada a uma capacidade funcional do sistema; usado nesta
dissertação como \emph{proxy} para ``requisito''. \\

\emph{Grafo de Requisitos} & Grafo direcionado derivado do
\emph{Call Graph} através de um mapeamento \code{func\_to\_req}, em
que cada vértice é um requisito de domínio e as arestas
representam dependências entre requisitos. \\

\emph{HITS} & Algoritmo de \citet{kleinberg1999} que atribui a cada
vértice dois \emph{scores}: \emph{hub} (orquestrador) e
\emph{authority} (dependência central). \\

\emph{PageRank} & Algoritmo de ranqueamento que modela a
probabilidade estacionária de um navegador aleatório visitar cada
vértice do grafo \parencite{page1999, brin1998}. \\

Kendall $\tau_b$ & Coeficiente de correlação de ranks baseado em
pares concordantes e discordantes, com correção para empates
\parencite{kendall1938, goodman1954}. Toma valores em $[-1, 1]$. \\

\emph{Rank Correlation} & Medida da semelhança entre duas
ordenações; nesta dissertação é utilizada para comparar rankings
produzidos por diferentes algoritmos. \\

\emph{Random Walk} & Processo estocástico que descreve a
trajetória de um caminhante que se move aleatoriamente entre
vértices de um grafo, seguindo probabilidades associadas às
arestas. \\

\emph{SALSA} & Algoritmo de \citet{lempel2001} que combina HITS
com cadeias de Markov sobre um grafo bipartido, normalizando
localmente por grau. \\

Spearman $\rho$ & Coeficiente de correlação de Pearson aplicado a
ranks; mede a força e direção da relação monotónica entre duas
ordenações \parencite{spearman1904}. Toma valores em $[-1, 1]$. \\

\emph{Threats to validity} & Conjunto de ameaças à
generalização e à interpretação de um estudo empírico, organizadas
por \citet{wohlin2012} em quatro categorias:
\emph{construct}, \emph{internal}, \emph{external}, \emph{conclusion}. \\

\bottomrule
\end{longtable}
